% Irodov, Arihant PDF pp. 31--33, printed pp. 22--24.
\exproblem{131}{22}
某球对称力场中的势能为
$U(r)=a/r^2-b/r$，其中 $a,b>0$。
求：（a）平衡半径及其稳定性；
（b）最大吸引力，并说明 $U(r)$、径向力 $F_r(r)$ 的图像特点。
\exsolution
径向力为
\[
F_r(r)=-\frac{\mathrm dU}{\mathrm dr}
=\frac{2a}{r^3}-\frac{b}{r^2}.
\]
平衡需 $F_r=0$，故
$\boxed{r_0=2a/b}$。
在此点
\[
U''(r_0)=\left(\frac{6a}{r^4}
-\frac{2b}{r^3}\right)_{r=r_0}
=\frac{2a}{r_0^4}>0,
\]
势能极小，平衡稳定。
当 $r>r_0$ 时 $F_r<0$，吸引力大小
$f(r)=b/r^2-2a/r^3$。
求导 $f'=-2b/r^3+6a/r^4$，
最大值出现在 $r_m=3a/b$，有
\[
\boxed{f_{\max}=\frac{b^3}{27a^2}}.
\]
\excheck $r\to0^+$ 时 $U\to+\infty,F_r\to+\infty$
（强斥力）；$r\to\infty$ 时 $U\to0^-$、$F_r\to0^-$。
$U$ 在 $r_0$ 取负极小，$F_r$ 在 $r_0$ 过零，
在 $r_m$ 取最小负值。

\exproblem{132}{22}
二维势场的势能为
$U(x,y)=\alpha x^2+\beta y^2$，
其中 $\alpha,\beta>0$ 且 $\alpha\ne\beta$。
求：（a）是否为中心力场；
（b）等势线和力大小相同的点构成什么曲线，
它们的半轴比各是多少。
\exsolution
力矢量
\[
\bm F=-\nabla U=(-2\alpha x,-2\beta y).
\]
若是以原点为中心的中心力，任一点的 $\bm F$
必须与 $\bm r=(x,y)$ 平行。
但 $\bm r\times\bm F$
的平面垂直分量为 $2(\alpha-\beta)xy$，
一般不为零，因此
$\boxed{\text{非中心力场}}$。
对于给定势能 $U_0>0$，等势线
$\alpha x^2+\beta y^2=U_0$ 是椭圆；
其 $x,y$ 轴半径分别为
$\sqrt{U_0/\alpha},\sqrt{U_0/\beta}$，故
$\boxed{a_x/a_y=\sqrt{\beta/\alpha}}$。
力大小相同 $|\bm F|=F_0$ 则
\[
\alpha^2x^2+\beta^2y^2=\frac{F_0^2}{4},
\quad\boxed{a_x/a_y=\frac{\beta}{\alpha}}.
\]
\excheck 两族曲线都是同轴椭圆，但长短轴比不同；
仅 $\alpha=\beta$ 时力才处处径向。

\exproblem{133}{22}
考察两种与时间无关的平面力场：
$\bm F_1=ay\,\bm i$ 和
$\bm F_2=ax\,\bm i+by\,\bm j$，
其中 $a,b$ 均为常数。
求两力场分别是否存在势能函数。
\exsolution
在单连通平面上，足够光滑力场存在势函数的必要充分条件是
$\partial F_y/\partial x=\partial F_x/\partial y$。
对 $\bm F_1$，
\[
\frac{\partial F_{1y}}{\partial x}=0,\qquad
\frac{\partial F_{1x}}{\partial y}=a.
\]
因此当 $a\ne0$ 时不满足条件，
$\boxed{\bm F_1\text{ 非保守}}$；
$a=0$ 时它为零场，当然保守。
对 $\bm F_2$，两偏导均为零，故保守，
相应势能可取
\[
\boxed{U_2(x,y)=-\frac12ax^2-\frac12by^2+C},
\]
因为 $-\nabla U_2=\bm F_2$。
\excheck 即使力与时间无关也未必是保守力；
判断标准是旋度而不是“定常”。

\exproblem{134}{23}
质量为 $m$ 的物块以初速 $v_0$
沿倾角 $\alpha$ 的粗糙斜面向上滑动，
摩擦系数为 $k$。
求停下以前的路程及摩擦力在此过程中所做的功。
\exsolution
沿斜面向上运动时，重力切向分量及摩擦力均指向下方，
减速度大小
$a=g(\sin\alpha+k\cos\alpha)$。
由 $0=v_0^2-2as$：
\[
\boxed{s_{\rm stop}=
\frac{v_0^2}{2g(\sin\alpha+k\cos\alpha)}}.
\]
摩擦力大小 $kmg\cos\alpha$，恒反向于位移，
所以
\[
W_f=-kmg\cos\alpha\,s_{\rm stop}
=\boxed{-\frac{mv_0^2 k}
 {2(\tan\alpha+k)}}.
\]
\excheck $k=0$ 时停下距离为 $v_0^2/(2g\sin\alpha)$
且摩擦功为零；$k>0$ 时摩擦功为负。

\exproblem{135}{23}
图1.30中，小圆片从高度 $H$ 的光滑坡顶由静止滑下，
坡道末端有一个高于地面 $h$ 的水平段，
圆片从该段末端水平飞出。
求使落地点与平台边缘的水平距离 $s$ 最大的 $h$
及相应最大射程。
\exsolution
能量守恒给出水平飞出速度
$v=\sqrt{2g(H-h)}$。
从高 $h$ 的平台水平抛出，落地用时
$t=\sqrt{2h/g}$，故
\[
s(h)=vt=2\sqrt{h(H-h)}.
\]
平方为
$s^2=4[hH-h^2]=H^2-4(h-H/2)^2$，
所以在 $0<h<H$ 中，唯一最大点为
\[
\boxed{h_{\rm opt}=\frac H2,\qquad
s_{\max}=H}.
\]
\excheck $h\to0$ 或 $h\to H$ 时射程均趋零，
验证内部极大值合理。

\exproblem{136}{23}
图1.31中，小物体从高度 $h$ 由静止出发，
沿光滑槽道下滑，槽道随后变为
半径 $R=h/2$ 的竖直半圆弧。
物体在接近顶部前离开槽道，之后作抛体运动。
求离槽后自由飞行过程中到达\emph{最高点}时的速率。
\exsolution
取半圆弧的圆心距地面高度为 $R=h/2$，
用 $\theta$ 表示圆弧上质点半径与竖直向上的夹角。
其高度 $y=R(1+\cos\theta)$。
由能量守恒
\[
v^2=2g(h-y)=2gR(1-\cos\theta).
\]
在槽内向心方向列力学方程
$N+mg\cos\theta=mv^2/R$；
当 $N=0$ 时离开槽道，所以
$g\cos\theta=2g(1-\cos\theta)$，
即 $\cos\theta=2/3$、$v^2=2gR/3=gh/3$。
离开时速度沿圆弧切线，其水平方向分量为
$v_x=v\cos\theta$；之后忽略阻力，
到抛体最高点的速率即 $|v_x|$，故
\[
\boxed{v_{\rm apex}=\frac23\sqrt{\frac{gh}{3}}}.
\]
\excheck 此“最高点”是离槽后的\emph{自由抛物线顶点}，
不是圆槽最上端；二者不能混淆。

\exproblem{137}{23}
摆长 $\ell$ 的单摆起初竖直静止。
将悬点突然以某水平方向的恒速移动，
使小球相对悬点获得初速。
求让小球能够在整个竖直圆周内绳始终绷紧所需的
悬点最小速率，以及小球运动至水平位置时的绳张力。
\exsolution
设悬点移动速率为 $u$；
悬点突然运动时小球速度尚未来得及变化，
故小球相对悬点的初速大小为 $u$，
初始位置在圆周最低点。
从最低点到最高点的重力势能增量为 $2mg\ell$，
能量守恒给出
$v_{\rm top}^2=u^2-4g\ell$。
为了在圆顶保持绳绷紧，
由向心方程 $T_{\rm top}+mg=mv_{\rm top}^2/\ell$
要求 $v_{\rm top}^2\ge g\ell$。
故
\[
\boxed{u_{\min}=\sqrt{5g\ell}}.
\]
在水平位置高度较起点上升 $\ell$，
此刻速度平方
$v_{\rm side}^2=u_{\min}^2-2g\ell=3g\ell$。
重力沿指向悬点的径向分量为零，
于是
\[
\boxed{T_{\rm side}=\frac{mv_{\rm side}^2}{\ell}=3mg}.
\]
\excheck 临界最小初速对应最高点张力恰为零，
再小便会在到达最高点之前绳子松弛。

\exproblem{138}{23}
如图1.32，从上方看，一质量很小的圆片 $A$
用长为 $\ell_0$ 的水平细绳系于静止圆柱
表面切点 $B$，圆柱半径为 $R$。
初始圆片获得垂直于细绳的速度 $v_0$，
随运动细绳不断绕圆柱缠绕。
忽略摩擦，求圆片最终撞上圆柱的时间。
\exsolution
细绳张力始终垂直于圆片的瞬时速度，不做功，
所以圆片速率保持 $v_0$。
令绳子在柱面上新缠绕的圆心角为 $\varphi$，
则自由绳长
$\ell(\varphi)=\ell_0-R\varphi$。
圆片轨迹是圆柱横截面的渐开线，其弧长微元
\[
\mathrm ds=\ell(\varphi)\,\mathrm d\varphi
=(\ell_0-R\varphi)\,\mathrm d\varphi.
\]
从 $\varphi=0$ 到绳全绕完
$\varphi_*=\ell_0/R$，总路程为
\[
s=\int_0^{\ell_0/R}
(\ell_0-R\varphi)\,\mathrm d\varphi
=\frac{\ell_0^2}{2R}.
\]
故
\[
\boxed{t=\frac{s}{v_0}=\frac{\ell_0^2}{2Rv_0}}.
\]
\excheck 绳长变短时圆片的瞬时曲率半径也逐渐变小；
不能把它的运动当作以固定半径 $\ell_0$ 的圆周运动。

\exproblem{139}{23}
图1.33中，长为 $\ell$、劲度系数为 $k$ 的轻弹性绳
上端固定在 $O$，下端带轻挡块 $B$。
质量为 $m$ 的小环 $A$ 从 $O$ 处静止开始下落，
先自由落下距离 $\ell$，随后挡块受力使弹性绳伸长。
求绳子的最大伸长量 $\Delta\ell$。
\exsolution
从初始位置到最大伸长点，
小环总下降距离为 $\ell+\Delta\ell$，
初末速率均为零。
忽略绳和挡块的质量以及耗散，
重力势能损失全部转为弹性势能：
\[
mg(\ell+\Delta\ell)=\frac12 k(\Delta\ell)^2.
\]
这是关于 $\Delta\ell$ 的一元二次方程，
取正根得
\[
\boxed{\Delta\ell=
\frac{mg}{k}\left(1+
\sqrt{1+\frac{2k\ell}{mg}}\right)}.
\]
\excheck $\ell\to0$ 时退化为从自然长位置
自由下落的竖直弹簧最大伸长 $2mg/k$；
另一负根没有物理意义。

\exproblem{140}{24}
图1.34中，质量为 $m$ 的物块 $A$
在光滑水平面上，初始由左侧绳 $PA$ 固定，
向右的绳跨过轻滑轮连接等质量重物 $B$。
另有弹簧竖直连接 $A$ 与其正上方固定点 $O$，
弹簧自然长 $\ell_0=0.50\,\mathrm m$，
劲度系数 $k=5mg/\ell_0$。
烧断左绳后，求 $A$ 刚开始离开水平面时的速度。
\exsolution
设 $A$ 向右滑过 $x$，由右绳长度守恒，
$B$ 下降同样距离 $x$，两物块速率相同为 $v$。
弹簧长度变为
$L=\sqrt{\ell_0^2+x^2}$。
弹簧向上的作用力分量为
$k(L-\ell_0)\ell_0/L$；
离面临界时地面支持力 $N=0$，故
\[
k(L-\ell_0)\frac{\ell_0}{L}=mg.
\]
代入 $k=5mg/\ell_0$ 得
$5(1-\ell_0/L)=1$，
所以 $L=5\ell_0/4$，进而
$x=3\ell_0/4$。
系统重物的势能损失为 $mgx$，
弹簧势能增加为
$\tfrac12 k(L-\ell_0)^2=5mg\ell_0/32$。
两等质量物块总动能为
$2(\tfrac12 mv^2)=mv^2$，故
\[
mv^2=mg\frac{3\ell_0}{4}
-\frac{5mg\ell_0}{32}
=\frac{19mg\ell_0}{32}.
\]
因此
\[
\boxed{v=\sqrt{\frac{19g\ell_0}{32}}
\simeq1.70\,\mathrm{m/s}}.
\]
\excheck 此刻弹簧竖直分力恰好抵消重力，
而非弹簧长度刚超过自然长时立刻离面。
